% Options for packages loaded elsewhere
\PassOptionsToPackage{unicode}{hyperref}
\PassOptionsToPackage{hyphens}{url}
\documentclass[
  12pt,
]{article}
\usepackage{xcolor}
\usepackage[margin=1in]{geometry}
\usepackage{amsmath,amssymb}
\setcounter{secnumdepth}{-\maxdimen} % remove section numbering
\usepackage{iftex}
\ifPDFTeX
  \usepackage[T1]{fontenc}
  \usepackage[utf8]{inputenc}
  \usepackage{textcomp} % provide euro and other symbols
\else % if luatex or xetex
  \usepackage{unicode-math} % this also loads fontspec
  \defaultfontfeatures{Scale=MatchLowercase}
  \defaultfontfeatures[\rmfamily]{Ligatures=TeX,Scale=1}
\fi
\usepackage{lmodern}
\ifPDFTeX\else
  % xetex/luatex font selection
  \setmonofont[]{Menlo}
\fi
% Use upquote if available, for straight quotes in verbatim environments
\IfFileExists{upquote.sty}{\usepackage{upquote}}{}
\IfFileExists{microtype.sty}{% use microtype if available
  \usepackage[]{microtype}
  \UseMicrotypeSet[protrusion]{basicmath} % disable protrusion for tt fonts
}{}
\makeatletter
\@ifundefined{KOMAClassName}{% if non-KOMA class
  \IfFileExists{parskip.sty}{%
    \usepackage{parskip}
  }{% else
    \setlength{\parindent}{0pt}
    \setlength{\parskip}{6pt plus 2pt minus 1pt}}
}{% if KOMA class
  \KOMAoptions{parskip=half}}
\makeatother
\usepackage{longtable,booktabs,array}
\newcounter{none} % for unnumbered tables
\usepackage{calc} % for calculating minipage widths
% Correct order of tables after \paragraph or \subparagraph
\usepackage{etoolbox}
\makeatletter
\patchcmd\longtable{\par}{\if@noskipsec\mbox{}\fi\par}{}{}
\makeatother
% Allow footnotes in longtable head/foot
\IfFileExists{footnotehyper.sty}{\usepackage{footnotehyper}}{\usepackage{footnote}}
\makesavenoteenv{longtable}
% definitions for citeproc citations
\NewDocumentCommand\citeproctext{}{}
\NewDocumentCommand\citeproc{mm}{%
  \begingroup\def\citeproctext{#2}\cite{#1}\endgroup}
\makeatletter
 % allow citations to break across lines
 \let\@cite@ofmt\@firstofone
 % avoid brackets around text for \cite:
 \def\@biblabel#1{}
 \def\@cite#1#2{{#1\if@tempswa , #2\fi}}
\makeatother
\newlength{\cslhangindent}
\setlength{\cslhangindent}{1.5em}
\newlength{\csllabelwidth}
\setlength{\csllabelwidth}{3em}
\newenvironment{CSLReferences}[2] % #1 hanging-indent, #2 entry-spacing
 {\begin{list}{}{%
  \setlength{\itemindent}{0pt}
  \setlength{\leftmargin}{0pt}
  \setlength{\parsep}{0pt}
  % turn on hanging indent if param 1 is 1
  \ifodd #1
   \setlength{\leftmargin}{\cslhangindent}
   \setlength{\itemindent}{-1\cslhangindent}
  \fi
  % set entry spacing
  \setlength{\itemsep}{#2\baselineskip}}}
 {\end{list}}
\usepackage{calc}
\newcommand{\CSLBlock}[1]{\hfill\break\parbox[t]{\linewidth}{\strut\ignorespaces#1\strut}}
\newcommand{\CSLLeftMargin}[1]{\parbox[t]{\csllabelwidth}{\strut#1\strut}}
\newcommand{\CSLRightInline}[1]{\parbox[t]{\linewidth - \csllabelwidth}{\strut#1\strut}}
\newcommand{\CSLIndent}[1]{\hspace{\cslhangindent}#1}
\setlength{\emergencystretch}{3em} % prevent overfull lines
\providecommand{\tightlist}{%
  \setlength{\itemsep}{0pt}\setlength{\parskip}{0pt}}
\usepackage{bookmark}
\IfFileExists{xurl.sty}{\usepackage{xurl}}{} % add URL line breaks if available
\urlstyle{same}
\hypersetup{
  pdftitle={librla: A library of randomized linear algebra routines},
  hidelinks,
  pdfcreator={LaTeX via pandoc}}

\title{librla: A library of randomized linear algebra routines}
\author{Adrianna Gillman\textsuperscript{1} \quad Zydrunas Gimbutas\textsuperscript{2} \\[6pt]
\small \textsuperscript{1}University of Colorado, Boulder, Department of Applied Mathematics \\
\small \textsuperscript{2}National Institute of Standards and Technology (NIST)}
\date{April 23, 2026}

\begin{document}
\maketitle

\section{Summary}\label{summary}

Randomized linear algebra algorithms have become a vital tool for a
variety of areas including fast direct solvers, reduced order modeling,
and data science. Additionally, randomized linear algebra provides
useful routines for solving total least squares problems and
rank-deficient least squares problems, doing matrix approximation, and
skeletonizing (i.e.~subset selection) a matrix (Chan and Hansen 1992).
\texttt{librla} is designed for small to mid-range sized matrices
(i.e.~up to roughly 10,000 in size depending on computing resources).
\texttt{librla} is not intended for matrices that are larger or do not
fit in memory.

\texttt{librla} includes the following options, all of which can be used
for both real and complex matrices:

\begin{itemize}
\tightlist
\item
  \texttt{qr\_sketch}: Randomized QR factorization.
\item
  \texttt{svd\_sketch}: Randomized Singular Value Decomposition (SVD).
\item
  \texttt{id\_sketch}: Interpolatory Decomposition via randomized
  sampling.
\end{itemize}

The user has the choice of specifying a desired accuracy (tolerance) or
rank. For each method, the user has the option to use extra samples and
to use a specified number of power iterations.\\
The package does include the option to create low-rank factorizations of
matrices that are applied via matrix-vector multiplication codes. When
the matrix-vector multiplication is ``matrix-free,'' the creation of the
low factorization is also matrix free. In order to use this option,
\texttt{librla} requires the ability to apply both the matrix and its
transpose via a subroutine.

\section{Statement of need}\label{statement-of-need}

While there is a large amount of research activity in the field of
randomized linear algebra, an easy-to-use, efficient and stable
factorization library is not available. For several factorization
libraries, the stability issue manifest itself by either failure in the
form of an invalid factorization or being unable to return anything.
When new techniques implemented in a high level language are dependent
on low level language libraries, it has a big impact on the portability
and usability of the new work. For example, this problem has impacted of
the field of fast direct solvers for boundary integral equations which
often depend on a legacy Fortran interpolatory decomposition library
(Martinsson et al. 2008).

Currently, there are two related routines in widely available,
maintained Python packages. They are PyTorch's \texttt{svd\_lowrank} and
SciPy's \texttt{id\_decomp} found in the
\texttt{scipy.linalg.interpolative} library. Codes allowing the user to
compare the performance are available in the \emph{compare} folder of
our repository. To illustrate the performance of librla, a subset of the
provided examples are presented here. The experiments were run on a
MacBook Pro with an M4 processor and 64 GB of RAM in the virtual
environment created by \texttt{setup\_venv.sh}. The experiments consider
the following matrices: Hilbert matrices, a matrix where the spectrum
decays \(\sim1/k\) for integer \(k\) between 1 and the size of the
matrix, and a matrix from a Gaussian mixture model (Dong et al. 2025).
Table 1 reports a subset of results obtained by running
\texttt{compare\_svd\_torch.py}. For a sequence of matrices, rank 15
approximate singular value decompositions are computed. This comparison
is only for the fixed rank factorizations as that is the only option
available in the Pytorch \texttt{svd\_lowrank} software. Additionally,
the singular value decomposition is the only factorization available in
Pytorch. The results demonstrate that the timings are comparable. Table
2 reports a subset of the result obtained by running
\texttt{compare\_id\_scipy.py}. The Interpolatory decomposition is the
only factorization available in SciPy and it allows the user to input a
desire rank or tolerance. The results in the table include both types of
experiments. Table 2 also reports the ratio of time it takes librla to
compute the factorization over the time it takes SciPy under the title
\emph{Speedup}. The results illustrate the benefits of using the
\texttt{librla} package.

\begin{longtable}[]{@{}
  >{\raggedright\arraybackslash}p{(\linewidth - 8\tabcolsep) * \real{0.3421}}
  >{\raggedright\arraybackslash}p{(\linewidth - 8\tabcolsep) * \real{0.1711}}
  >{\raggedright\arraybackslash}p{(\linewidth - 8\tabcolsep) * \real{0.1447}}
  >{\raggedright\arraybackslash}p{(\linewidth - 8\tabcolsep) * \real{0.1316}}
  >{\raggedright\arraybackslash}p{(\linewidth - 8\tabcolsep) * \real{0.2105}}@{}}
\caption{The table reports the relative error and the time in seconds
for creating rank 15 factorizations of different matrices using Pytorch
and librla.}\tabularnewline
\toprule\noalign{}
\begin{minipage}[b]{\linewidth}\raggedright
Problem
\end{minipage} & \begin{minipage}[b]{\linewidth}\raggedright
Size
\end{minipage} & \begin{minipage}[b]{\linewidth}\raggedright
Pytorch
\end{minipage} & \begin{minipage}[b]{\linewidth}\raggedright
librla
\end{minipage} & \begin{minipage}[b]{\linewidth}\raggedright
Relative Error
\end{minipage} \\
\midrule\noalign{}
\endfirsthead
\toprule\noalign{}
\begin{minipage}[b]{\linewidth}\raggedright
Problem
\end{minipage} & \begin{minipage}[b]{\linewidth}\raggedright
Size
\end{minipage} & \begin{minipage}[b]{\linewidth}\raggedright
Pytorch
\end{minipage} & \begin{minipage}[b]{\linewidth}\raggedright
librla
\end{minipage} & \begin{minipage}[b]{\linewidth}\raggedright
Relative Error
\end{minipage} \\
\midrule\noalign{}
\endhead
\bottomrule\noalign{}
\endlastfoot
Hilbert matrix & 2000 x 1000 & 0.002s & 0.0018s & 3.679e-08 \\
Hilbert matrix & 4000 x 2000 & 0.0045s & 0.0043s & 1.604e-07 \\
Decaying spectrum matrix & 800 x 600 & 0.0026s & 0.0023s & 1.529e-01 \\
Gaussian mixture model & 400 x 400 & 0.0016s & 0.0014s & 7.096e-01 \\
\end{longtable}

\begin{longtable}[]{@{}
  >{\raggedright\arraybackslash}p{(\linewidth - 12\tabcolsep) * \real{0.2549}}
  >{\raggedright\arraybackslash}p{(\linewidth - 12\tabcolsep) * \real{0.1275}}
  >{\raggedright\arraybackslash}p{(\linewidth - 12\tabcolsep) * \real{0.1961}}
  >{\raggedright\arraybackslash}p{(\linewidth - 12\tabcolsep) * \real{0.0882}}
  >{\raggedright\arraybackslash}p{(\linewidth - 12\tabcolsep) * \real{0.0882}}
  >{\raggedright\arraybackslash}p{(\linewidth - 12\tabcolsep) * \real{0.0882}}
  >{\raggedright\arraybackslash}p{(\linewidth - 12\tabcolsep) * \real{0.1569}}@{}}
\caption{The table reports the relative error and the time in seconds
for creating low factorizations of different matrices using SciPy and
librla. The speedup over the SciPy package is also reported. The
examples involve fixed rank or tolerance approximations.}\tabularnewline
\toprule\noalign{}
\begin{minipage}[b]{\linewidth}\raggedright
Problem
\end{minipage} & \begin{minipage}[b]{\linewidth}\raggedright
Size
\end{minipage} & \begin{minipage}[b]{\linewidth}\raggedright
Stopping Condition
\end{minipage} & \begin{minipage}[b]{\linewidth}\raggedright
SciPy
\end{minipage} & \begin{minipage}[b]{\linewidth}\raggedright
librla
\end{minipage} & \begin{minipage}[b]{\linewidth}\raggedright
Speedup
\end{minipage} & \begin{minipage}[b]{\linewidth}\raggedright
Relative Error
\end{minipage} \\
\midrule\noalign{}
\endfirsthead
\toprule\noalign{}
\begin{minipage}[b]{\linewidth}\raggedright
Problem
\end{minipage} & \begin{minipage}[b]{\linewidth}\raggedright
Size
\end{minipage} & \begin{minipage}[b]{\linewidth}\raggedright
Stopping Condition
\end{minipage} & \begin{minipage}[b]{\linewidth}\raggedright
SciPy
\end{minipage} & \begin{minipage}[b]{\linewidth}\raggedright
librla
\end{minipage} & \begin{minipage}[b]{\linewidth}\raggedright
Speedup
\end{minipage} & \begin{minipage}[b]{\linewidth}\raggedright
Relative Error
\end{minipage} \\
\midrule\noalign{}
\endhead
\bottomrule\noalign{}
\endlastfoot
Hilbert matrix & 2000 x 1000 & rank 15 & 0.0355s & 0.0024s & 14.6 &
8.571e-08 \\
Hilbert matrix & 4000 x 2000 & rank 15 & 0.1557s & 0.0046s & 33.9 &
1.306e-06 \\
Decaying spectrum matrix & 800 x 600 & tol = 0.01 & 0.2933s & 0.0146s &
20 & 0 \\
Gaussian mixture model & 400 x 400 & tol = 0.01 & 0.0699s & 0.0055s &
12.8 & 5.659e-05 \\
\end{longtable}

\section{State of the field}\label{state-of-the-field}

While the algorithm used in \texttt{librla} is built in a similar manner
to the randomized methods in (Halko et al. 2011; Liberty et al. 2007).
There is a large collection of related work. (Martinsson and Tropp 2020;
Mahoney 2011; Drineas and Mahoney 2017; Woodruff 2014; Voronin and
Martinsson 2017) are some review papers on the field. Some literature
(Mahoney and Drineas 2009; Sorensen and Embree 2016; Drineas et al.
2008) has focused on the design of CUR factorizations. While other
manuscripts (Frieze et al. 2004; Drineas, Kannan, et al. 2006b; Drineas,
Mahoney, et al. 2006) have focused on Monte Carlo approaches to creating
the low rank factorizations. Randomized projection methods (Halko et al.
2011; Liberty et al. 2007; Duersch and Gu 2020; Meier and Nakatsukasa
2024; Gu and Eisenstat 1996; Erichson et al. 2018; Rokhlin and Tygert
2008) are common. Deterministic factorization techniques exist (Duersch
and Gu 2017; Feng et al. 2018). Many techniques are specifically
designed for distributed memory and/or to not require many passes
through the data (Martinsson et al. 2019; Sarlos 2006; Tropp et al.
2017; Yu et al. 2017; Lindquist et al. 2020; Yang et al. 2015; Bjarkason
2019). There is much literature from the randomized linear algebra
community for applications including \(L_2\) solutions (Drineas, Kannan,
et al. 2006a; Drineas et al. 2007; Avron et al. 2010; Meng et al. 2014;
Pilanci and Wainwright 2016), reduced order modeling (Sorensen and
Embree 2016), linear programming (Drineas et al. 2012; Chowdhury et al.
2020), statistics (Drineas and Mahoney 2016), machine learning (Yao et
al. 2021), Newton methods (Roosta-Khorasani and Mahoney 2019; Pilanci
and Wainwright 2017), and differentially private matrices (Balu and
Furon 2016; Cormode et al. 2018; Choi et al. 2020). The user-specified
rank algorithm in \texttt{librla} is influenced by the
\texttt{svd\_lowrank} routine in PyTorch.

\section{Mathematics}\label{mathematics}

The algorithm that serves as the foundation of \texttt{librla} is called
\texttt{orth\_sketch}. It creates an orthogonal basis of the range of
the operator of interest. The user-specified tolerance portion of the
package was heavily influenced by the methods presented in (Halko et al.
2011; Martinsson et al. 2008; Liberty et al. 2007). The user-specified
rank portion of the package was influenced by the \texttt{svd\_lowrank}
in PyTorch. Once the orthogonal basis is created, it is possible to
create the low-rank QR, SVD or interpolatory decomposition via standard
techniques. Pseudocodes for all of the factorizations are provided in
the \texttt{PSEUDOCODE.md} file found in the \emph{distrib} folder.

Users call the desired factorization subroutine with a linear operator
\({\bf A}\) of size \(m\times n\), a desired rank and stopping condition
\({tol\_or\_rank}\). The user has the option to use power iteration and
specify the number of iterations \(power\_iter\) to accelerate
convergence. Users are also provided the option to specify a number of
oversampling vectors. Roughly speaking, given a linear operator
\({\bf A}\), tolerance condition or desired rank \(tol\_or\_rank\),
specified \(block\_size\) and \(power\_iter\), \texttt{orth\_sketch}
randomly samples the range of \({\bf A}\) to create an orthogonal basis
for the range. The range is sampled by applying \({\bf A}\) to a matrix
\({ \Omega }\) of size \(n \times block\_size\) whose entries are
uniformly randomly sampled from \([-1,1]\). This choice of matrix
entries for \(\Omega\) ensures that the vectors are linearly independent
but may not be the optimal basis choice. (This is an open question.) The
size \(\Omega\) is set so that number of columns includes the over
sampling vectors; i.e.~the parameter \(block\_size\) is updated to be
the sum of the original \(block\_size\) and the number of oversampling
vectors. In the tolerance option, the \(block\_size\) is adaptively
increased until the stopping criterion is satisfied. In the rank option,
the size of \(\Omega\) is determined by the desired rank and the number
of user specified extra samples. The orthogonal basis is formed by
taking a pivoted QR factorization of \({\bf A}\Omega\), where \(\Omega\)
is optionally refined via power iteration with reorthogonalization at
each step. The magnitudes of the diagonal entries of \({\bf R}\) are
stored in a vector \({\bf diagR}\) and are used in the stopping
criterion for the tolerance option of the factorization.
\texttt{orth\_sketch} returns \({\bf Q}\) or a submatrix of \({\bf Q}\),
\({\bf diagR}\) and an error flag to the factorization routine that
called it. The matrix \({\bf Q}\) is then used to create the desired
factorization. In the tolerance option of \texttt{librla}, the range of
\({\bf A}\) is considered sufficiently if the ratio of the magnitude of
the largest and smallest diagonal entries of the \({\bf R}\) is less
than the specified tolerance \({tol\_or\_rank}\). If the sketching step
terminates early, the factorization routine falls back to a
deterministic algorithm.

\section{Software design}\label{software-design}

The library \texttt{librla}, written natively in Python, Julia and
Matlab, provides low-rank QR factorizations, SVDs and interpolatory
decompositions. The algorithms randomly sample the range of the matrix
or operator in a similar manner to (Halko et al. 2011). A key feature of
this package is that it is designed to exploit Level 3 BLAS operators as
much as possible. The use of Level 3 BLAS allows for all the linear
algebraic operations in the method be executed via low level optimized
code. This design choice is the foundation for the efficiency of the
software. In fact, the performance in terms of both speed and stability
is the same in all three languages. The codes are consistent between
languages allowing for the software to be easily maintained.

\section{Research impact statement}\label{research-impact-statement}

The avialability of efficient and stable randomized factorization
techniques has the potential to impact various research areas. For
example, the \texttt{librla} library allows for the fast direct solvers
the authors develop to be easily shared with collaborators.
Additionally, the \texttt{librla} library allows for easy
reproducibility of the solver work as it is no longer dependent on
platform dependent wrappers on FORTRAN code. Other areas that can
experience similar benefits include reduced order modeling and data
compression.

\section{AI usage disclosure}\label{ai-usage-disclosure}

Generative AI tools were used during the preparation of this work to
assist with code development, testing, and editing of the manuscript.
The authors reviewed and verified all AI-generated content and take full
responsibility for the accuracy and integrity of the final publication.
The original algorithm was prototyped and debugged in Python. Claude
Code (Anthropic, Claude Sonnet 3.5 and Claude Opus 4.6) was used to
assist in porting to Matlab and Julia, while maintaining consistency in
API calls and documentation. Unit tests were generated to ensure
correctness of the ported code.

\section{Acknowledgements}\label{acknowledgements}

The work by A. Gillman was supported by the National Science Foundation
(DMS-2110886), and a Knut and Alice Wallenberg Foundation Grant. Part of
this work was carried out while A. Gillman was in residence at Institut
Mittag-Leffler in Djursholm, Sweden in autumn 2025, supported by the
Swedish Research Council under grant no. 2021-06594.

Contributions by staff of NIST, an agency of the U.S. Government, are
not subject to copyright within the United States.

Certain commercial software and equipment are identified in this paper
to foster understanding. Such identification does not imply
recommendation or endorsement by NIST, nor does it imply that the
software or equipment identified is necessarily the best available for
the purpose.

\section*{References}\label{references}
\addcontentsline{toc}{section}{References}

\protect\phantomsection\label{refs}
\begin{CSLReferences}{1}{1}
\bibitem[\citeproctext]{ref-Avron:2010}
Avron, Haim, Petar Maymounkov, and Sivan Toledo. 2010. {``Blendenpik:
Supercharging LAPACK's Least-Squares Solver.''} \emph{SIAM Journal on
Scientific Computing} 32 (3): 1217--36.
\url{https://doi.org/10.1137/090767911}.

\bibitem[\citeproctext]{ref-Balu:2016}
Balu, Raghavendran, and Teddy Furon. 2016. {``Differentially Private
Matrix Factorization Using Sketching Techniques.''} \emph{Proceedings of
the 4th ACM Workshop on Information Hiding and Multimedia Security} (New
York, NY, USA). \url{https://doi.org/10.1145/2909827.2930793}.

\bibitem[\citeproctext]{ref-Bjarkason:2019}
Bjarkason, Elvar K. 2019. {``Pass-Efficient Randomized Algorithms for
Low-Rank Matrix Approximation Using Any Number of Views.''} \emph{SIAM
Journal on Scientific Computing} 41 (4): A2355--83.
\url{https://doi.org/10.1137/18M118966X}.

\bibitem[\citeproctext]{ref-Chan:1992}
Chan, Tony F., and Per Christian Hansen. 1992. {``Some Applications of
the Rank Revealing QR Factorization.''} \emph{SIAM Journal on Scientific
and Statistical Computing} 13 (3): 727--41.
\url{https://doi.org/10.1137/0913043}.

\bibitem[\citeproctext]{ref-Choi:2020}
Choi, Seung Geol, Dana Dachman-Soled, Mukul Kulkarni, and Arkady
Yerukhimovich. 2020. {``Differentially-Private Multi-Party Sketching for
Large-Scale Statistics.''} \emph{Proceedings on Privacy Enhancing
Technologies} 3: 153--74.
\url{https://doi.org/10.2478/popets-2020-0047}.

\bibitem[\citeproctext]{ref-Chowdhury:2020}
Chowdhury, Agniva, Palma London, Haim Avron, and Petros Drineas. 2020.
{``Faster Randomized Infeasible Interior Point Methods for Tall/Wide
Linear Programs.''} \emph{Proceedings of the 34th International
Conference on Neural Information Processing Systems} (Red Hook, NY,
USA), NIPS '20.

\bibitem[\citeproctext]{ref-Cormode:2018}
Cormode, Graham, Somesh Jha, Tejas Kulkarni, Ninghui Li, Divesh
Srivastava, and Tianhao Wang. 2018. {``Privacy at Scale: Local
Differential Privacy in Practice.''} \emph{Proceedings of the 2018
International Conference on Management of Data} (New York, NY, USA).
\url{https://doi.org/10.1145/3183713.3197390}.

\bibitem[\citeproctext]{ref-Dong:2025}
Dong, Yijun, Chao Chen, Per-Gunnar Martinsson, and Katherine Pearce.
2025. {``Robust Blockwise Random Pivoting: Fast and Accurate Adaptive
Interpolative Decomposition.''} \emph{SIAM Journal on Matrix Analysis
and Applications} 46 (3): 1791--815.
\url{https://doi.org/10.1137/24M1678027}.

\bibitem[\citeproctext]{ref-Drineas:2006}
Drineas, Petros, Ravi Kannan, and Michael W. Mahoney. 2006a. {``Fast
Monte Carlo Algorithms for Matrices i: Approximating Matrix
Multiplication.''} \emph{SIAM Journal on Computing} 36 (1): 132--57.
\url{https://doi.org/10.1137/S0097539704442684}.

\bibitem[\citeproctext]{ref-Drineas:2006_2}
Drineas, Petros, Ravi Kannan, and Michael W. Mahoney. 2006b. {``Fast
Monte Carlo Algorithms for Matrices II: Computing a Low-Rank
Approximation to a Matrix.''} \emph{SIAM Journal on Computing} 36 (1):
158--83. \url{https://doi.org/10.1137/S0097539704442696}.

\bibitem[\citeproctext]{ref-Drineas:2012}
Drineas, Petros, Malik Magdon-Ismail, Michael W. Mahoney, and David P.
Woodruff. 2012. {``Fast Approximation of Matrix Coherence and
Statistical Leverage.''} \emph{Journal of Maching Learning Research} 13
(1): 3475--506.

\bibitem[\citeproctext]{ref-Drineas:2016}
Drineas, Petros, and Michael W. Mahoney. 2016. \emph{RandNLA: Randomized
Numerical Linear Algebra}. (New York, NY, USA) 59 (6).
\url{https://doi.org/10.1145/2842602}.

\bibitem[\citeproctext]{ref-Drineas2017LecturesOR}
Drineas, Petros, and Michael W. Mahoney. 2017. {``Lectures on Randomized
Numerical Linear Algebra.''} \emph{ArXiv} abs/1712.08880.
\url{https://doi.org/10.1090/PCMS/025/01}.

\bibitem[\citeproctext]{ref-Drinea:2006_3}
Drineas, Petros, Michael W. Mahoney, and S. Muthukrishnan. 2006.
{``Sampling Algorithms for L2 Regression and Applications.''}
\emph{Proceedings of the Seventeenth Annual ACM-SIAM Symposium on
Discrete Algorithm}, SODA '06, 1127--36.

\bibitem[\citeproctext]{ref-Drineas:2008}
Drineas, Petros, Michael W. Mahoney, and S. Muthukrishnan. 2008.
{``Relative-Error \$CUR\$ Matrix Decompositions.''} \emph{SIAM Journal
on Matrix Analysis and Applications} 30 (2): 844--81.
\url{https://doi.org/10.1137/07070471X}.

\bibitem[\citeproctext]{ref-Drineas:2007}
Drineas, Petros, Michael W. Mahoney, S. Muthukrishnan, and Tamás Sarlós.
2007. {``Faster Least Squares Approximation.''} \emph{Numerische
Mathematik} 117: 219--49.
\url{https://doi.org/10.1007/s00211-010-0331-6}.

\bibitem[\citeproctext]{ref-Duersch:2017_2}
Duersch, Jed A., and Ming Gu. 2017. {``Randomized QR with Column
Pivoting.''} \emph{SIAM Journal on Scientific Computing} 39 (4):
C263--91. \url{https://doi.org/10.1137/15M1044680}.

\bibitem[\citeproctext]{ref-Duersch:2020}
Duersch, Jed A., and Ming Gu. 2020. {``Randomized Projection for
Rank-Revealing Matrix Factorizations and Low-Rank Approximations.''}
\emph{SIAM Review} 62 (3): 661--82.
\url{https://doi.org/10.1137/20M1335571}.

\bibitem[\citeproctext]{ref-ERICHSON:2018}
Erichson, N. Benjamin, Ariana Mendible, Sophie Wihlborn, and J. Nathan
Kutz. 2018. {``Randomized Nonnegative Matrix Factorization.''}
\emph{Pattern Recognition Letters} 104: 1--7.
\url{https://doi.org/10.1016/j.patrec.2018.01.007}.

\bibitem[\citeproctext]{ref-Feng:2018}
Feng, Yuehua, Jianwei Xiao, and Ming Gu. 2018. {``Randomized Complete
Pivoting for Solving Symmetric Indefinite Linear Systems.''} \emph{SIAM
Journal on Matrix Analysis and Applications} 39 (4): 1616--41.
\url{https://doi.org/10.1137/17M1150013}.

\bibitem[\citeproctext]{ref-Frieze:2004}
Frieze, Alan, Ravi Kannan, and Santosh Vempala. 2004. \emph{Fast
Monte-Carlo Algorithms for Finding Low-Rank Approximations}. (New York,
NY, USA) 51 (6). \url{https://doi.org/10.1145/1039488.1039494}.

\bibitem[\citeproctext]{ref-Gu:1996}
Gu, Ming, and Stanley C. Eisenstat. 1996. {``Efficient Algorithms for
Computing a Strong Rank-Revealing QR Factorization.''} \emph{SIAM
Journal on Scientific Computing} 17 (4): 848--69.
\url{https://doi.org/10.1137/0917055}.

\bibitem[\citeproctext]{ref-Halko:2011}
Halko, N., P. G. Martinsson, and J. A. Tropp. 2011. {``Finding Structure
with Randomness: Probabilistic Algorithms for Constructing Approximate
Matrix Decompositions.''} \emph{SIAM Review} 53 (2): 217--88.
\url{https://doi.org/10.1137/090771806}.

\bibitem[\citeproctext]{ref-Liberty:2007}
Liberty, Edo, F. Woolfe, P. G Martinsson, V. Rokhlin, and V. Tygert.
2007. {``Randomized Algorithms for the Low-Rank Approximation of
Matrices.''} \emph{Proceedings of the National Academy of Sciences} 104
(51): 20167--72. \url{https://doi.org/10.1073/pnas.0709640104}.

\bibitem[\citeproctext]{ref-Lindquist:2020}
Lindquist, Neil, Piotr Luszczek, and Jack Dongarra. 2020. {``{Replacing
Pivoting in Distributed Gaussian Elimination with Randomized Techniques
}.''} \emph{2020 IEEE/ACM 11th Workshop on Latest Advances in Scalable
Algorithms for Large-Scale Systems (ScalA)}, 35--43.
\url{https://doi.org/10.1109/ScalA51936.2020.00010}.

\bibitem[\citeproctext]{ref-Mahoney:2011}
Mahoney, Michael W. 2011. {``Randomized Algorithms for Matrices and
Data.''} \emph{Foundations and Trends in Machine Learning} 3 (2):
123--224. \url{https://doi.org/10.1561/2200000035}.

\bibitem[\citeproctext]{ref-Mahoney:2009}
Mahoney, Michael W., and Petros Drineas. 2009. {``CUR Matrix
Decompositions for Improved Data Analysis.''} \emph{Proceedings of the
National Academy of Sciences} 106 (3): 697--702.
\url{https://doi.org/10.1073/pnas.0803205106}.

\bibitem[\citeproctext]{ref-Martinsson:2019}
Martinsson, P. G., G. Quintana-Ortı́, and N. Heavner. 2019. {``randUTV: A
Blocked Randomized Algorithm for Computing a Rank-Revealing UTV
Factorization.''} \emph{ACM Trans. Math. Softw.} 45 (1).
\url{https://doi.org/10.1145/3242670}.

\bibitem[\citeproctext]{ref-Tygert:2008}
Martinsson, P. G., V. Rokhlin, Y. Shkolnisky, and M. Tygert. 2008.
\emph{{ID}: A Software Package for Low-Rank Approximation of Matrices
via Interpolative Decompositions}.
\url{https://tygert.com/software.html}.

\bibitem[\citeproctext]{ref-Martinsson:2020}
Martinsson, P.-G., and J. A. Tropp. 2020. {``Randomized Numerical Linear
Algebra: Foundations and Algorithms.''} \emph{Acta Numerica} 29:
403--572. \url{https://doi.org/10.1017/S0962492920000021}.

\bibitem[\citeproctext]{ref-MEIER:2024}
Meier, Maike, and Yuji Nakatsukasa. 2024. {``Fast Randomized Numerical
Rank Estimation for Numerically Low-Rank Matrices.''} \emph{Linear
Algebra and Its Applications} 686: 1--32.
\url{https://doi.org/10.1016/j.laa.2024.01.001}.

\bibitem[\citeproctext]{ref-Meng:2014}
Meng, Xiangrui, Michael A. Saunders, and Michael W. Mahoney. 2014.
{``LSRN: A Parallel Iterative Solver for Strongly over- or
Underdetermined Systems.''} \emph{SIAM Journal on Scientific Computing}
36 (2): C95--118. \url{https://doi.org/10.1137/120866580}.

\bibitem[\citeproctext]{ref-Pilanci:2016}
Pilanci, Mert, and Martin J Wainwright. 2016. {``Iterative Hessian
Sketch: Fast and Accurate Solution Approximation for Constrained
Least-Squares.''} \emph{Journal of Machine Learning Research} 17 (1):
1842--79.

\bibitem[\citeproctext]{ref-Pilanci:2017}
Pilanci, Mert, and Martin J. Wainwright. 2017. {``Newton Sketch: A Near
Linear-Time Optimization Algorithm with Linear-Quadratic Convergence.''}
\emph{SIAM Journal on Optimization} 27 (1): 205--45.
\url{https://doi.org/10.1137/15M1021106}.

\bibitem[\citeproctext]{ref-Rokhlin:2008}
Rokhlin, Vladimir, and Mark Tygert. 2008. {``A Fast Randomized Algorithm
for Overdetermined Linear Least-Squares Regression.''} \emph{Proceedings
of the National Academy of Sciences} 105 (36): 13212--17.
\url{https://doi.org/10.1073/pnas.0804869105}.

\bibitem[\citeproctext]{ref-Roosta-Khorasani:2019}
Roosta-Khorasani, Farbod, and Michael W. Mahoney. 2019.
\emph{Sub-Sampled Newton Methods}. (Berlin, Heidelberg) 174 (1--2).
\url{https://doi.org/10.1007/s10107-018-1346-5}.

\bibitem[\citeproctext]{ref-Sarlos:2006}
Sarlos, Tamas. 2006. {``Improved Approximation Algorithms for Large
Matrices via Random Projections.''} \emph{Proceedings of the 47th Annual
IEEE Symposium on Foundations of Computer Science}, FOCS '06, 143--52.
\url{https://doi.org/10.1109/FOCS.2006.37}.

\bibitem[\citeproctext]{ref-Sorensen:2016}
Sorensen, D. C., and Mark Embree. 2016. {``A DEIM Induced CUR
Factorization.''} \emph{SIAM Journal on Scientific Computing} 38 (3):
A1454--82. \url{https://doi.org/10.1137/140978430}.

\bibitem[\citeproctext]{ref-Tropp:2017}
Tropp, Joel A., Alp Yurtsever, Madeleine Udell, and Volkan Cevher. 2017.
{``Practical Sketching Algorithms for Low-Rank Matrix Approximation.''}
\emph{SIAM Journal on Matrix Analysis and Applications} 38 (4):
1454--85. \url{https://doi.org/10.1137/17M1111590}.

\bibitem[\citeproctext]{ref-Voronin:2017}
Voronin, S., and P. G. Martinsson. 2017. {``Efficient Algorithms for Cur
and Interpolative Matrix Decompositions.''} \emph{Advances in
Computational Mathematics} 43 (3): 495--516.
\url{https://doi.org/10.1007/s10444-016-9494-8}.

\bibitem[\citeproctext]{ref-Woodruff:2014}
Woodruff, David P. 2014. \emph{Sketching as a Tool for Numerical Linear
Algebra}. (Hanover, MA, USA) 10 (1--2).
\url{https://doi.org/10.1561/0400000060}.

\bibitem[\citeproctext]{ref-Yang:2015}
Yang, Jiyan, Xiangrui Meng, and Michael W. Mahoney. 2015.
{``Implementing Randomized Matrix Algorithms in Parallel and Distributed
Environments.''} \emph{Proceedings of the IEEE} 104: 58--92.
\url{https://doi.org/10.1109/JPROC.2015.2494219}.

\bibitem[\citeproctext]{ref-Yao:2021}
Yao, Zhewei, Amir Gholami, Sheng Shen, Mustafa Mustafa, Kurt Keutzer,
and Michael Mahoney. 2021. {``ADAHESSIAN: An Adaptive Second Order
Optimizer for Machine Learning.''} \emph{Proceedings of the AAAI
Conference on Artificial Intelligence} 35 (12): 10665--73.
\url{https://doi.org/10.1609/aaai.v35i12.17275}.

\bibitem[\citeproctext]{ref-Yu:2017}
Yu, Wenjian, Yu Gu, Jian Li, Shenghua Liu, and Yaohang Li. 2017.
{``Single-Pass PCA of Large High-Dimensional Data.''} \emph{Proceedings
of the Twenty-Sixth International Joint Conference on Artificial
Intelligence, {IJCAI-17}}, 3350--56.
\url{https://doi.org/10.24963/ijcai.2017/468}.

\end{CSLReferences}

\end{document}
